%!TEX program = xelatex
%!TEX root = ../all_beamer.tex

% ============================================================
%  Nonparametric Regression -- Beamer Presentation
%  Style modelled on the SLT beamer (ml-slt-beamer.pdf)
% ============================================================

%\documentclass[aspectratio=169]{beamer}
%
%\usetheme{default}
%\usecolortheme{default}
%
%% ---- Colours --------------------------------------------
%\definecolor{HeaderBlue}{RGB}{35,55,85}
%\definecolor{AlertRed}{RGB}{140,30,30}
%\definecolor{AlertRed}{RGB}{240,220,220}
%\definecolor{GreenBlock}{RGB}{40,110,60}
%\definecolor{GreenLightBlock}{RGB}{210,235,215}
%\definecolor{TitleGold}{RGB}{220,180,60}
%
%% ---- Beamer colours -------------------------------------
%\setbeamercolor{frametitle}{bg=HeaderBlue,fg=white}
%\setbeamercolor{title}{fg=white}
%\setbeamercolor{author}{fg=white}
%\setbeamercolor{date}{fg=white}
%\setbeamercolor{institute}{fg=white}
%\setbeamercolor{structure}{fg=HeaderBlue}
%\setbeamercolor{section in head/foot}{bg=HeaderBlue,fg=white}
%\setbeamercolor{footline}{bg=HeaderBlue,fg=white}
%
%% ---- Frame title ----------------------------------------
%\setbeamertemplate{frametitle}{%
%  \nointerlineskip
%  \begin{beamercolorbox}[wd=\paperwidth,ht=2.4ex,dp=1.2ex,leftskip=0.5cm]{frametitle}%
%    \usebeamerfont{frametitle}\insertframetitle
%    \ifx\insertframesubtitle\empty\else
%      \quad{\color{TitleGold}\small\textit{\insertframesubtitle}}%
%    \fi
%  \end{beamercolorbox}%
%}
%
%% ---- Navigation bar -------------------------------------
%\setbeamertemplate{headline}{%
%  \begin{beamercolorbox}[wd=\paperwidth,ht=2.5ex,dp=1.125ex]{section in head/foot}%
%    \insertsectionnavigationhorizontal{\paperwidth}{}{}%
%  \end{beamercolorbox}%
%}
%
%% ---- Footline -------------------------------------------
%\setbeamertemplate{footline}{%
%  \begin{beamercolorbox}[wd=\paperwidth,ht=2.25ex,dp=1ex,
%      leftskip=0.5cm,rightskip=0.5cm]{footline}%
%    \footnotesize
%    Giorgio Gambosi (Tor Vergata)
%    \hfill Nonparametric Regression \hfill
%    A.Y.\ 2025--2026 \quad \insertframenumber/\inserttotalframenumber
%  \end{beamercolorbox}%
%}
%
%\setbeamertemplate{navigation symbols}{}
%\setbeamertemplate{itemize items}[circle]
%\setbeamertemplate{enumerate items}[circle]
%
%% ---- Fonts ----------------------------------------------
%\usepackage[T1]{fontenc}
%\usepackage[utf8]{inputenc}
%\usefonttheme{professionalfonts}
%% !TEX root = all_beamer.tex
% ==============================================================================
% HEADER LATEX PER PRESENTAZIONI BEAMER
% Corso di Machine Learning - Giorgio Gambosi
% Università di Roma "Tor Vergata"
% ==============================================================================

% ------------------------------------------------------------------------------
% CLASSE DOCUMENTO E METADATI
% ------------------------------------------------------------------------------
\documentclass[aspectratio=169, 10pt]{beamer}

\subtitle{Course of Machine Learning}
\author{Giorgio Gambosi}
\institute[Tor Vergata]{Master's Degree in Computer Science\\University of Rome ``Tor Vergata''}
\date{A.Y. 2025--2026}

% ------------------------------------------------------------------------------
% PACCHETTI MATEMATICI
% ------------------------------------------------------------------------------
\usepackage{amsmath, amssymb, amsfonts, mathtools, amsthm}
\usepackage{bm}              % Simboli matematici in grassetto
\usepackage{upgreek}         % Lettere greche dritte
\usepackage{derivative}      % Gestione semplificata delle derivate

% ------------------------------------------------------------------------------
% PACCHETTI GRAFICI E VISUALIZZAZIONE
% ------------------------------------------------------------------------------
\usepackage{graphicx}        % Inclusione immagini
\usepackage{xcolor}          % Gestione colori
\usepackage{xkcdcolors}      % Colori aggiuntivi stile xkcd
\usepackage{microtype}       % Miglioramenti tipografici

% ------------------------------------------------------------------------------
% PACCHETTI PER TABELLE
% ------------------------------------------------------------------------------
\usepackage{booktabs}        % Tabelle professionali
\usepackage{array}           % Estensioni per tabelle
\usepackage{multirow}        % Celle multi-riga
\usepackage{multicol}        % Layout multi-colonna

% ------------------------------------------------------------------------------
% TIKZ E PGFPLOTS
% ------------------------------------------------------------------------------
\usepackage{tikz}
\usepackage{tikz-network}
\usepackage[customcolors]{hf-tikz}
\usepackage{pgfplots}
\pgfplotsset{compat=1.18}
\usepgfplotslibrary{fillbetween}

% Librerie TikZ
\usetikzlibrary{arrows, arrows.meta, decorations.pathmorphing, decorations.pathreplacing}
\usetikzlibrary{decorations.markings, snakes, positioning, backgrounds, fit, shadows}
\usetikzlibrary{automata, patterns, calc, intersections, through, shapes}
\usetikzlibrary{shapes.geometric, shapes.arrows, scopes, chains, matrix}

% ------------------------------------------------------------------------------
% ALTRI PACCHETTI
% ------------------------------------------------------------------------------
\usepackage[skins]{tcolorbox}  % Box colorati
\usepackage{subcaption}         % Sottocaption per figure
\usepackage{listings}           % Inclusione codice sorgente

% ------------------------------------------------------------------------------
% TEMA E STILE BEAMER
% ------------------------------------------------------------------------------
\usetheme{Madrid}
\usecolortheme{whale}
\useinnertheme{rounded}
\usefonttheme{professionalfonts}
\setbeamertemplate{navigation symbols}{}  % Rimuove simboli di navigazione

% ------------------------------------------------------------------------------
% DEFINIZIONE COLORI PERSONALIZZATI
% ------------------------------------------------------------------------------
% Colori principali del tema
\definecolor{MainBlue}{RGB}{31,56,100}      % Blu principale (header, titoli)
\definecolor{AccentBlue}{RGB}{70,130,180}   % Blu accento (blocchi)
\definecolor{AlertRed}{RGB}{160,30,30}      % Rosso per alert
\definecolor{ExGreen}{RGB}{30,120,60}       % Verde per esempi
\definecolor{BlockBg}{RGB}{220,232,246}     % Sfondo blocchi
\definecolor{torgray}{RGB}{220,225,235}     % Grigio Tor Vergata

% Colori alternativi (Tor Vergata style)
\definecolor{TvDark}{RGB}{30,50,80}         % Blu scuro header
\definecolor{TvMid}{RGB}{70,100,140}        % Blu medio blocchi
\definecolor{TvLight}{RGB}{210,220,235}     % Blu chiaro sfondi
\definecolor{TvAlert}{RGB}{160,30,30}       % Rosso alert
\definecolor{TvGreen}{RGB}{20,100,50}       % Verde esempi

% Colori supplementari
\definecolor{torred}{RGB}{153,0,0}
\definecolor{torcyan}{RGB}{0,112,192}
\definecolor{cerise}{HTML}{DA3A62}
\definecolor{teal}{HTML}{029386}
\definecolor{slateblue}{HTML}{5B5EA6}
\definecolor{orange}{HTML}{F97306}

% Colori per nodi TikZ
\definecolor{nodeblue}{RGB}{101,140,196}    % xkcd Faded Blue

% ------------------------------------------------------------------------------
% CONFIGURAZIONE COLORI BEAMER
% ------------------------------------------------------------------------------
\setbeamercolor{frametitle}{bg=MainBlue, fg=white}
\setbeamercolor{title}{bg=MainBlue, fg=white}
\setbeamercolor{structure}{fg=MainBlue}

% Blocchi standard
\setbeamercolor{block title}{bg=AccentBlue, fg=white}
\setbeamercolor{block body}{bg=BlockBg}

% Blocchi alert
\setbeamercolor{block title alerted}{bg=AlertRed, fg=white}
\setbeamercolor{block body alerted}{bg=red!8}

% Blocchi esempio
\setbeamercolor{block title example}{bg=ExGreen, fg=white}
\setbeamercolor{block body example}{bg=green!7}

% ------------------------------------------------------------------------------
% CONFIGURAZIONE FONT BEAMER
% ------------------------------------------------------------------------------
\setbeamerfont{frametitle}{size=\normalsize,series=\bfseries}
\setbeamerfont{block title}{size=\small,series=\bfseries}

% ------------------------------------------------------------------------------
% STILI TIKZ PERSONALIZZATI
% ------------------------------------------------------------------------------
\tikzset{
  >=latex,  % Frecce stile LaTeX di default
  % Nodi per reti neurali
  node/.style={thick,circle,draw=myblue,minimum size=22,inner sep=0.5,outer sep=0.6},
  node in/.style={node,green!20!black,draw=mygreen!30!black,fill=mygreen!25},
  node hidden/.style={node,blue!20!black,draw=myblue!30!black,fill=myblue!20},
  node convol/.style={node,orange!20!black,draw=myorange!30!black,fill=myorange!20},
  node out/.style={node,red!20!black,draw=myred!30!black,fill=myred!20},
  % Connessioni
  connect/.style={thick,mydarkblue},
  connect arrow/.style={-{Latex[length=4,width=3.5]},thick,mydarkblue,shorten <=0.5,shorten >=1},
  % Stili numerati per facile mappatura
  node 1/.style={node in},
  node 2/.style={node hidden},
  node 3/.style={node out}
}

% ------------------------------------------------------------------------------
% AMBIENTI TEOREMI
% ------------------------------------------------------------------------------
\setbeamertemplate{theorems}[numbered]
\newtheorem{mytheorem}{Theorem}
\newtheorem{mydefinition}{Definition}

% ------------------------------------------------------------------------------
% CONFIGURAZIONE LISTINGS (CODICE)
% ------------------------------------------------------------------------------
\lstset{
  language=Python,
  basicstyle=\ttfamily\tiny,
  keywordstyle=\color{MainBlue}\bfseries,
  commentstyle=\color{gray},
  frame=single,
  breaklines=true,
}

% ==============================================================================
% MACRO MATEMATICHE PERSONALIZZATE
% ==============================================================================

% ------------------------------------------------------------------------------
% SPAZI CALLIGRAFICI (Calligraphic spaces)
% ------------------------------------------------------------------------------
\providecommand{\calX}{\mathcal{X}}         % Spazio input
\providecommand{\calY}{\mathcal{Y}}         % Spazio output
\providecommand{\calH}{\mathcal{H}}         % Spazio ipotesi
\providecommand{\calT}{\mathcal{T}}         % Training set
\providecommand{\calA}{\mathcal{A}}         % Algoritmo
\providecommand{\calR}{\mathcal{R}}         % Rischio
\providecommand{\calF}{\mathcal{F}}         % Famiglia di funzioni
\providecommand{\calL}{\mathcal{L}}         % Loss function

% Alias per compatibilità
\providecommand{\Rr}{\mathcal{R}}
\providecommand{\Tr}{\mathcal{T}}
\providecommand{\Hr}{\mathcal{H}}

% ------------------------------------------------------------------------------
% RISCHIO E PROBABILITÀ
% ------------------------------------------------------------------------------
\providecommand{\barR}{\overline{\mathcal{R}}}  % Rischio empirico
\providecommand{\EmpR}{\overline{\mathcal{R}}}  % Rischio empirico (alias)
\providecommand{\Rbar}{\overline{\mathcal{R}}}  % Rischio empirico (alias)
\providecommand{\pM}{p_M}                       % Probabilità modello
\providecommand{\pC}{p_C}                       % Probabilità corretta
\providecommand{\Prob}[2]{\mathbb{P}_{#1}\!\left[#2\right]}  % Probabilità

% ------------------------------------------------------------------------------
% FUNZIONI INDICATRICI
% ------------------------------------------------------------------------------
\providecommand{\indicator}[1]{\mathbf{1}\!\left[#1\right]}  % Funzione indicatrice
\providecommand{\ind}[1]{\mathbf{1}[#1]}                      % Versione breve

% ------------------------------------------------------------------------------
% OPERATORI MATEMATICI
% ------------------------------------------------------------------------------
\providecommand{\argmax}{\operatorname*{arg\,max}}  % Argomento del massimo
\providecommand{\argmin}{\operatorname*{arg\,min}}  % Argomento del minimo
\providecommand{\vcd}[1]{\operatorname{VCdim}(#1)} % VC dimension
\providecommand{\defeq}{\stackrel{\text{def}}{=}}  % Definito come

% ------------------------------------------------------------------------------
% VETTORI E MATRICI (Bold Math)
% ------------------------------------------------------------------------------
% Vettori minuscoli
\providecommand{\bx}{\mathbf{x}}            % Vettore x
\providecommand{\bz}{\mathbf{z}}            % Vettore z
\providecommand{\bw}{\mathbf{w}}            % Vettore peso
\providecommand{\btt}{\mathbf{t}}           % Vettore target
\providecommand{\wvec}{\mathbf{w}}          % Vettore peso (alias)
\providecommand{\bth}{\boldsymbol{\theta}}  % Parametro theta

% Matrici maiuscole
\providecommand{\bX}{\mathbf{X}}            % Matrice dati
\providecommand{\bZ}{\mathbf{Z}}            % Matrice Z
\providecommand{\Xmat}{\mathbf{X}}          % Matrice dati (alias)
\providecommand{\Hmat}{\mathbf{H}}          % Matrice H
\providecommand{\Smat}{\mathbf{S}}          % Matrice scatter generica
\providecommand{\bS}{\mathbf{S}}            % Matrice scatter
\providecommand{\Ii}{\mathbf{I}}            % Matrice identità
\providecommand{\bZero}{\mathbf{0}}         % Vettore zero
\providecommand{\bPsi}{\boldsymbol{\Psi}}   % Matrice Psi

% Medie (overline)
\providecommand{\ox}{\overline{\mathbf{x}}} % Media vettore x
\providecommand{\ow}{\overline{\mathbf{w}}} % Media vettore w
\providecommand{\oX}{\overline{\mathbf{X}}} % Media matrice X

% Matrici scatter specifiche (analisi discriminante)
\providecommand{\SW}{\mathbf{S}_W}          % Within-class scatter
\providecommand{\SB}{\mathbf{S}_B}          % Between-class scatter
\providecommand{\ST}{\mathbf{S}_T}          % Total scatter
\providecommand{\GX}{G_{\mathbf{X}}}        % Matrice Gram

% ------------------------------------------------------------------------------
% COMANDI DI FORMATTAZIONE TESTO
% ------------------------------------------------------------------------------
\providecommand{\term}[1]{\textbf{\color{accent}#1}}        % Termine importante
\providecommand{\halert}[1]{{\color{AlertRed}\textbf{#1}}}  % Alert rosso
\providecommand{\mlalert}[1]{{\color{AccentBlue}\textbf{#1}}}  % Alert blu
\providecommand{\mlemph}[1]{{\color{MainBlue}\textit{#1}}}  % Enfasi blu

% ==============================================================================
% FINE HEADER
% ==============================================================================

\newcommand{\blankpage}{
\clearpage
\thispagestyle{empty}
\null
\clearpage
}

% ---------- Packages ----------
%\usepackage{amsmath, amssymb, mathtools}
%\usepackage{booktabs}
%\usepackage{tikz}
%\usetikzlibrary{arrows.meta, positioning, shapes.geometric}
%\usepackage{xcolor}
%\usepackage{microtype}
%\usepackage{graphicx}
%% indicator function via amsmath
%
%% ---------- Macros ----------
%
%\newcommand{\calX}{\mathcal{X}}
%\newcommand{\calY}{\mathcal{Y}}
%\newcommand{\calH}{\mathcal{H}}
%\newcommand{\calT}{\mathcal{T}}
%\newcommand{\calA}{\mathcal{A}}
%\newcommand{\calR}{\mathcal{R}}
%\newcommand{\barR}{\overline{\mathcal{R}}}
%\newcommand{\pM}{p_M}
%\newcommand{\pC}{p_C}
%\newcommand{\indicator}[1]{\mathbf{1}\!\left[#1\right]}
%\newcommand{\argmax}{\operatorname*{arg\,max}}
%\newcommand{\argmin}{\operatorname*{arg\,min}}

% ---------- Title ----------
\title[Nonparametric regression]{Nonparametric regression}
%\subtitle{Course of Machine Learning}
%\author{Giorgio Gambosi}
%\institute[Tor Vergata]{Master's Degree in Computer Science\\University of Rome ``Tor Vergata''}
%\date{A.Y. 2025–2026}

% ---- Math -----------------------------------------------
%\usepackage{amsmath,amsfonts,amssymb,mathtools}
%\usepackage{bm}
%\usepackage{upgreek}
%
%% ---- Graphics -------------------------------------------
%\usepackage{graphicx}
%
%% ---- tcolorbox ------------------------------------------
%\usepackage[skins]{tcolorbox}

%\newtcolorbox{block}[1]{%
%  enhanced, colback=BlockBg, colframe=MainBlue,
%  fonttitle=\bfseries\small, title=#1,
%  left=1pt, right=4pt, top=1pt, bottom=1pt}

%\newtcolorbox{alertblock}[1]{%
%  enhanced, colback=AlertRed, colframe=AlertRed,
%  fonttitle=\bfseries\small, title=#1,
%  left=4pt, right=4pt, top=2pt, bottom=2pt}

%\newtcolorbox{greenblock}[1]{%
%  enhanced, colback=GreenLightBlock, colframe=GreenBlock,
%  fonttitle=\bfseries\small, title=#1,
%  left=4pt, right=4pt, top=2pt, bottom=2pt}

% ---- Math macros (names that avoid LaTeX conflicts) -----

%\newcommand{\Ii}{\mathbf{I}}
%\newcommand{\bZero}{\mathbf{0}}
%\newcommand{\bX}{\mathbf{X}}
%\newcommand{\bZ}{\mathbf{Z}}
%\newcommand{\btt}{\mathbf{t}}
%\newcommand{\bS}{\mathbf{S}}
%\newcommand{\bx}{\mathbf{x}}
%\newcommand{\bz}{\mathbf{z}}
%\newcommand{\bw}{\mathbf{w}}
%\newcommand{\ox}{\overline{\mathbf{x}}}
%\newcommand{\ow}{\overline{\mathbf{w}}}
%\newcommand{\oX}{\overline{\mathbf{X}}}
%\newcommand{\bPsi}{\boldsymbol{\Psi}}
%\newcommand{\halert}[1]{{\color{AlertRed}\textbf{#1}}}

% ---- Title slide template --------------------------------
%\setbeamertemplate{title page}{%
%  \begin{beamercolorbox}[wd=\paperwidth,ht=\paperheight]{HeaderBlue}%
%    \vspace{2cm}
%    \begin{center}
%      {\color{white}\Large\bfseries\inserttitle}\\[0.6em]
%      {\color{TitleGold}\large\insertsubtitle}\\[1.8em]
%      {\color{white}\normalsize\insertauthor}\\[0.4em]
%      {\color{white}\small\insertinstitute}\\[1em]
%      {\color{white}\small\insertdate}
%    \end{center}
%  \end{beamercolorbox}%
%}
%
%\title{Nonparametric Regression}
%\subtitle{Course of Machine Learning}
%\author{Giorgio Gambosi}
%\institute{Master's Degree in Computer Science\\[2pt]%
%           University of Rome ``Tor Vergata''}
%\date{A.Y.\ 2025--2026}

% ============================================================
%\begin{document}

\begin{frame}[plain]
  \titlepage
\end{frame}

\begin{frame}{Outline}
  \tableofcontents
\end{frame}

% ============================================================
\section{From Bayesian Regression to Kernel Methods}
% ============================================================

%-------------------------------------------------------
\begin{frame}{Fully Bayesian Regression}{Recap}
  \begin{block}{Bayesian prediction framework}
    No single optimal parameter $\overline{\mathbf{w}}^*$ is identified.
    The \halert{predictive distribution} $p(t|\bx)$ is derived directly:
    \[
      p(t|\bx,\bX,\btt,\alpha,\beta)=\mathcal{N}\!\left(t\mid m(\bx),\,\sigma^2(\bx)\right)
    \]
    with
    \[
      m(\bx)=\beta\,\ox^T\bS\,\oX^T\btt,\qquad
      \sigma^2(\bx)=\tfrac{1}{\beta}+\ox^T\bS\,\ox
    \]
    and $\bS=(\alpha\Ii+\beta\oX^T\oX)^{-1}\in\Real^{(d+1)\times(d+1)}$.
  \end{block}

  \medskip
  \begin{itemize}
    \item Prediction is the expectation: $h(\bx)=m(\bx)$.
    \item No explicit parameter optimisation -- sum over all possible models.
  \end{itemize}
\end{frame}

%-------------------------------------------------------
\begin{frame}{Prediction as a Weighted Sum of Targets}{Deriving the equivalent kernel}
  \begin{itemize}
    \item Since $\oX^T\btt=\sum_{i=1}^n\ox_i t_i$, the prediction becomes:
      \[
        h(\bx)=\beta\,\ox^T\bS\sum_{i=1}^n\ox_i t_i
              =\sum_{i=1}^n\underbrace{\beta\,\ox^T\bS\ox_i}_{\kappa(\bx,\bx_i)}\,t_i
      \]
    \item Prediction is a \halert{weighted sum of all target values} $t_i$.
    \item Weights depend on both $\bx_i$ (training) and $\bx$ (query).
  \end{itemize}

  \medskip
  \begin{block}{Equivalent kernel}
    \[
      \kappa(\bx_1,\bx_2)=\beta\,\ox_1^T\bS\ox_2
      \;\Longrightarrow\;
      h(\bx)=\sum_{i=1}^n\kappa(\bx,\bx_i)\,t_i
    \]
    \begin{itemize}
      \item $\kappa(\bx_1,\bx_2)$ measures the \emph{dependence} between targets at $\bx_1$ and $\bx_2$.
      \item With basis functions $\Vphi$: $\kappa(\bx_1,\bx_2)=\beta\,\Vphi(\bx_1)^T\bS\,\Vphi(\bx_2)$.
    \end{itemize}
  \end{block}
\end{frame}

%-------------------------------------------------------
\begin{frame}{Equivalent Kernel -- Localisation Property}{}
  \begin{columns}
    \column{0.50\textwidth}
      \begin{itemize}
        \item $\kappa(\bx,\bx_i)$ gives higher weight to training points $\bx_i$ \halert{near} $\bx$.
        \item Natural localisation: distant points receive small weight.
        \item Left: $\kappa(x,x_i)$ vs.\ $x$ for three values of $x_i$.
        \item Right: surface plot on the $(x,x_i)$ plane.
      \end{itemize}
    \column{0.48\textwidth}
      \begin{center}
        \includegraphics[width=\linewidth]{figures/equivkernel.pdf}
      \end{center}
  \end{columns}
  \smallskip
  \begin{center}
    \small\textit{Equivalent kernel for Gaussian basis functions, one input feature.}
  \end{center}
\end{frame}

%-------------------------------------------------------
\begin{frame}{Equivalent Kernel -- Different Basis Functions}{}
  \begin{itemize}
    \item The localisation property holds for \emph{any} choice of basis functions $\Vphi$.
    \item Below: $\kappa(0,x)$ for polynomial (left) and Gaussian (right) basis functions.
  \end{itemize}
  \begin{center}
    \begin{tabular}{ccc}
      \includegraphics[width=0.38\linewidth]{figs/Figure3-11a.pdf} &
      \hspace{3mm} &
      \includegraphics[width=0.38\linewidth]{figs/Figure3-11b.pdf}
    \end{tabular}
  \end{center}
  \begin{itemize}
    \item Key insight: instead of designing $\Vphi$ first, define a \halert{localised kernel} directly
          on $\Real^d\times\Real^d$ and use it for prediction.
  \end{itemize}
\end{frame}

% ============================================================
\section{Kernel Regression}
% ============================================================

%-------------------------------------------------------
\begin{frame}{Kernel Regression}{Setup and goal}
  \begin{block}{Objective}
    Predict $t$ for any query $\bx$ by weighting training targets
    according to their \halert{proximity} to $\bx$, controlled by a kernel function.
  \end{block}

  \medskip
  \begin{itemize}
    \item A \halert{kernel function} $\kappa_h(\bx)$ returns non-negligible values only near $\bZero$.
    \item $h$ is a bandwidth hyperparameter controlling the locality.
    \item Regression target: approximate the conditional expectation
      \[
        \mathbb{E}[t\mid\bx]=
        \frac{\displaystyle\int t\,p(\bx,t)\,dt}{\displaystyle\int p(\bx,t)\,dt}
      \]
  \end{itemize}
\end{frame}

%-------------------------------------------------------
\begin{frame}{The Gaussian (RBF) Kernel}{}
  \begin{block}{Most common kernel choice}
    \[
      g(\bx)=e^{-\dfrac{\|\bx\|^2}{2h^2}}
    \]
    Plotted below for $d=1$ and several values of $h$.
  \end{block}

  \begin{center}
    \includegraphics[width=0.68\linewidth]{figures/gaussian_kernel.png}
  \end{center}
  \begin{itemize}
    \item Larger $h$ $\Rightarrow$ wider kernel, more smoothing.
    \item Smaller $h$ $\Rightarrow$ narrower kernel, more local prediction.
  \end{itemize}
\end{frame}

%-------------------------------------------------------
\begin{frame}{Nadaraya-Watson Estimator}{Derivation}
  \begin{itemize}
    \item Approximate the joint density:
      $\displaystyle p(\bx,t)\approx\tfrac{1}{n}\sum_{i=1}^n\kappa_h(\bx-\bx_i)\kappa_h(t-t_i)$.
    \item Assume $\kappa_h$ is non-negative, zero-mean, unit area (a density).
    \item Then $\int\kappa_h(t-t_i)\,dt=1$ and $\int t\,\kappa_h(t-t_i)\,dt=t_i$.
  \end{itemize}

  \medskip
  \begin{block}{Nadaraya-Watson estimator}
    \[
      h(\bx)=\frac{\displaystyle\sum_{i=1}^n\kappa_h(\bx-\bx_i)\,t_i}
                  {\displaystyle\sum_{i=1}^n\kappa_h(\bx-\bx_i)}
             =\sum_{i=1}^n w_i(\bx)\,t_i
    \]
    where $\displaystyle w_i(\bx)=\dfrac{\kappa_h(\bx-\bx_i)}{\sum_{j=1}^n\kappa_h(\bx-\bx_j)}$.
  \end{block}
\end{frame}

%-------------------------------------------------------
\begin{frame}{Nadaraya-Watson -- Properties and Extension}{}
  \begin{itemize}
    \item Prediction is a \halert{normalised linear combination} of all target values.
    \item Weights $w_i(\bx)$ sum to $1$ for every $\bx$.
    \item Nearby training points receive the highest weights.
  \end{itemize}

  \medskip
  \begin{block}{Extension with basis functions $\Vphi$}
    \[
      h(\bx)=\sum_{i=1}^n w_i(\Vphi(\bx))\,t_i,\qquad
      w_i(\bx)=\frac{\kappa_h(\Vphi(\bx)-\Vphi(\bx_i))}%
                    {\sum_{j=1}^n\kappa_h(\Vphi(\bx)-\Vphi(\bx_j))}
    \]
  \end{block}

  \medskip
  \begin{alertblock}{Limitation}
    Prediction is a weighted average of \emph{constants} (the $t_i$).
    More expressiveness is obtained by fitting a local linear model instead.
  \end{alertblock}
\end{frame}

% ============================================================
\section{Locally Weighted Regression}
% ============================================================

%-------------------------------------------------------
\begin{frame}{Locally Weighted Regression (LOESS)}{Motivation}
  \begin{block}{Improvement over Nadaraya-Watson}
    LOESS fits a \halert{local linear model} around each query point $\bx$,
    replacing the simple weighted average of constants.
  \end{block}

  \medskip
  \begin{itemize}
    \item Minimise a \halert{locally weighted} loss function.
    \item Weights $\psi_i(\bx)$ depend on the distance between $\bx$ and $\bx_i$.
    \item Assuming basis functions $\Vphi$ and $\ox_i=\Vphi(\bx_i)$:
  \end{itemize}

  \[
    L(\bx)=\sum_{i=1}^n\psi_i(\bx)\bigl(\ow^T\ox_i-t_i\bigr)^2,
    \qquad
    \psi_i(\bx)=\kappa_h(\bx-\bx_i)
  \]

  \begin{itemize}
    \item A \emph{different} weighted model is fitted for each query $\bx$.
  \end{itemize}
\end{frame}

%-------------------------------------------------------
\begin{frame}{Locally Weighted Regression -- Solution}{}
  \begin{block}{Weighted least squares closed-form solution}
    The minimisation
    \[
      \ow^*(\bx)=\arg\min_{\bw}\sum_{i=1}^n\psi_i(\bx)\bigl(\ow^T\ox_i-t_i\bigr)^2
    \]
    has the solution
    \[
      \ow^*(\bx)=\bigl(\oX^T\bPsi(\bx)\oX\bigr)^{-1}\oX^T\bPsi(\bx)\,\btt
    \]
    where $\bPsi(\bx)$ is the $n\times n$ diagonal matrix with $\bPsi(\bx)_{ii}=\psi_i(\bx)$.
  \end{block}

  \medskip
  \begin{itemize}
    \item Prediction: $h(\bx)=\ow^*(\bx)^T\ox$.
    \item A \halert{different linear model} is solved for each query point $\bx$.
    \item More expressive than Nadaraya-Watson.
    \item Higher computational cost: one linear system per prediction.
  \end{itemize}
\end{frame}

% ============================================================
\section{Local Logistic Regression}
% ============================================================

%-------------------------------------------------------
\begin{frame}{Local Logistic Regression}{Extension to classification}
  \begin{block}{Idea}
    Apply the same local-weighting philosophy to \halert{classification}
    via a logistic regression model.
  \end{block}

  \medskip
  \begin{itemize}
    \item Replace the squared loss with a \halert{weighted cross-entropy}:
      \[
        L(\bx)=\sum_{i=1}^n\kappa_h(\bx-\bx_i)
               \bigl[t_i\log p_i-(1-t_i)\log(1-p_i)\bigr]
      \]
      with $p_i=\sigma(\ow^T\ox_i)$ (logistic sigmoid).
    \item The loss is \emph{maximised} with respect to $\bw$.
    \item Weights $\kappa_h(\bx-\bx_i)$ emphasise training points close to $\bx$.
    \item Prediction: $h(\bx)=\sigma\bigl(\ow^*(\bx)^T\ox\bigr)$.
  \end{itemize}
\end{frame}

% ============================================================
\section{Gaussian Processes}
% ============================================================

%-------------------------------------------------------
\begin{frame}{Gaussian Processes}{Motivation}
  \begin{block}{Key idea}
    Perform inference \emph{directly in the space of functions}
    $f:\Real^d\to\Real$, using a \halert{Gaussian process} (GP)
    as a prior distribution over functions.
  \end{block}

  \medskip
  \begin{itemize}
    \item \textbf{Stochastic process}: collection of random variables
          $\{f(\bx):\bx\in\chi\}$ -- the values taken by $f$ on domain $\chi$.
    \item \textbf{Gaussian process}: a stochastic process such that for \emph{any} finite
          subset $\bX=(\bx_1,\ldots,\bx_n)\subset\chi$, the values
          $f(\bx_1),\ldots,f(\bx_n)$ have a \halert{joint multivariate Gaussian} distribution.
    \item GPs generalise finite-dimensional multivariate Gaussians to \emph{infinite} dimensions.
    \item Common case: $\chi=\Real^d$; finite domains also allowed.
  \end{itemize}
\end{frame}

%-------------------------------------------------------
\begin{frame}{Specifying a Gaussian Process}{}
  A GP is fully specified by two functions defined on $\chi$:

  \medskip
  \begin{block}{Mean function $m:\chi\to\Real$}
    $m(\bx_i)=\mathbb{E}[f(\bx_i)]$.
    Mean vector: $\Vmu_{\bX}=(m(\bx_1),\ldots,m(\bx_n))^T$.
    Often set to $\bZero$ without loss of generality.
  \end{block}

  \medskip
  \begin{block}{Covariance (kernel) function $\kappa:\chi\times\chi\to\Real$}
    Required to be a \halert{positive definite kernel}.
    Covariance matrix: $\VSigma_{\bX}[i,j]=\kappa(\bx_i,\bx_j)$.
    The kernel encodes assumptions about smoothness, length scale, etc.
  \end{block}

  \medskip
  Notation: $f(\bx)\sim\mathcal{GP}(m(\bx),\,\kappa(\bx,\bx'))$.
\end{frame}

%-------------------------------------------------------
\begin{frame}{Positive Definite Kernels}{}
  \begin{block}{Positive definiteness}
    $\kappa$ is positive definite if, for any $n$ distinct points
    $\bx_1,\ldots,\bx_n$ and constants $c_1,\ldots,c_n$ (not all zero):
    \[
      \sum_{i=1}^n\sum_{j=1}^n c_i c_j\,\kappa(\bx_i,\bx_j)>0
    \]
  \end{block}

  \medskip
  \begin{itemize}
    \item Equivalently: the \halert{Gram matrix} $G_{\bX}$ has all positive eigenvalues:
      \[
        G_{\bX}=\begin{pmatrix}
          \kappa(\bx_1,\bx_1) & \cdots & \kappa(\bx_1,\bx_n)\\
          \vdots & \ddots & \vdots\\
          \kappa(\bx_n,\bx_1) & \cdots & \kappa(\bx_n,\bx_n)
        \end{pmatrix}
      \]
    \item Large collection of valid kernels known; composable by standard rules.
    \item Kernel shape determines the \emph{character} of sampled functions
          (smoothness, range, periodicity, \ldots).
  \end{itemize}
\end{frame}

%-------------------------------------------------------
\begin{frame}{Gaussian Process -- Formal Summary}{}
  \begin{block}{GP distribution for any finite set}
    Given $\mathcal{GP}(m,\kappa)$, for any finite $\bX=(\bx_1,\ldots,\bx_n)$:
    \[
      p(f(\bx_1),\ldots,f(\bx_n))=\mathcal{N}(f;\,\Vmu_{\bX},\,\VSigma_{\bX})
    \]
    where $\Vmu_{\bX}[i]=m(\bx_i)$ and $\VSigma_{\bX}[i,j]=\kappa(\bx_i,\bx_j)=G_{\bX}[i,j]$.
  \end{block}

  \medskip
  \begin{itemize}
    \item Any marginal of a GP distribution is Gaussian (by definition and marginalization).
    \item Sampling: draw from $\mathcal{N}(\Vmu_{\bX},G_{\bX})$ to obtain function values at $\bX$.
    \item The GP is a \halert{distribution over functions}: the kernel encodes their shape.
    \item Different kernels produce different types of sampled functions.
  \end{itemize}
\end{frame}

% ============================================================
\section{Kernels and Prior Samples}
% ============================================================

%-------------------------------------------------------
\begin{frame}{The RBF (Gaussian) Kernel}{}
  \begin{block}{Radial Basis Function kernel}
    \[
      \kappa(\bx_1,\bx_2)=\sigma^2\,e^{-\dfrac{\|\bx_1-\bx_2\|^2}{2\tau^2}}
    \]
    \begin{itemize}
      \item Covariance is high when $\bx_1$ and $\bx_2$ are close.
      \item $\tau$: length scale.  Larger $\tau$ $\Rightarrow$ smoother functions.
      \item Most widely used kernel in GP practice.
    \end{itemize}
  \end{block}

  \medskip
  \begin{itemize}
    \item Sampled functions tend to be \halert{smooth}: nearby inputs yield similar outputs.
    \item Degree of smoothness grows with $\tau$.
  \end{itemize}
\end{frame}

%-------------------------------------------------------
\begin{frame}{Prior Samples -- Large Length Scale $\tau$}{}
  \begin{center}
    \includegraphics[width=0.85\linewidth]{figures/rbf1.jpg}
  \end{center}
  \begin{itemize}
    \item RBF kernel with \halert{larger $\tau$}: samples are very smooth.
    \item Function values change slowly across the domain.
  \end{itemize}
\end{frame}

%-------------------------------------------------------
\begin{frame}{Prior Samples -- Small Length Scale $\tau$}{}
  \begin{center}
    \includegraphics[width=0.85\linewidth]{figures/rbf4.jpg}
  \end{center}
  \begin{itemize}
    \item RBF kernel with \halert{smaller $\tau$}: samples oscillate more rapidly.
    \item Correlation drops quickly with distance (strong locality).
  \end{itemize}
\end{frame}

% ============================================================
\section{GP Posterior -- No Noise}
% ============================================================

%-------------------------------------------------------
\begin{frame}{GP as a Bayesian Prior}{}
  \begin{block}{Bayesian perspective}
    The GP $\mathcal{GP}(m(\bx),\kappa(\bx,\bx'))$ is a \halert{prior} $p(f)$ over functions,
    defined independently of any observed data.
  \end{block}

  \medskip
  \begin{itemize}
    \item Prior over function values at training points:
      $p(f)=\mathcal{N}(f;\,\Vmu_{\bX},\,\VSigma_{\bX})$.
    \item Observing $(\bx_i,t_i)$ updates the prior to a \halert{posterior}.
    \item Two main settings:
      \begin{enumerate}
        \item Noise-free observations: $t_i=f(\bx_i)$ exactly.
        \item Noisy observations: $t_i=f(\bx_i)+\varepsilon$.
      \end{enumerate}
  \end{itemize}
\end{frame}

%-------------------------------------------------------
\begin{frame}{Prior Samples Before Observations}{}
  \begin{center}
    \includegraphics[width=0.83\linewidth]{figures/gp_no_noise.png}
  \end{center}
  \begin{itemize}
    \item \textcolor{red}{Red}: unknown target function $f(x)$ to be approximated.
    \item \textcolor{blue}{Blue}: functions sampled from the GP prior $p(f)$ (no data observed yet).
  \end{itemize}
\end{frame}

%-------------------------------------------------------
\begin{frame}{Joint Distribution of Training and Test Values}{No noise}
  \begin{itemize}
    \item Assume exact observations: $t_i=f(\bx_i)$.
    \item Training: $\bX=(\bx_1,\ldots,\bx_m)$, $\btt=(t_1,\ldots,t_m)^T$.
    \item Test points: $\bZ=(\bz_1,\ldots,\bz_r)$.
  \end{itemize}

  \medskip
  \begin{block}{Joint distribution (no noise)}
    $(f(\bX),f(\bZ))$ is an $(m+r)$-dimensional Gaussian with:
    \[
      \VSigma_{(\bX,\bZ)}=
      \begin{pmatrix}G_{\bX} & G_{\bZ,\bX}^T \\ G_{\bZ,\bX} & G_{\bZ}\end{pmatrix},
      \qquad
      G_{\bZ,\bX}[i,j]=\kappa(\bz_i,\bx_j)
    \]
  \end{block}
\end{frame}

%-------------------------------------------------------
\begin{frame}{Gaussian Conditional Distributions}{Useful property for GP inference}
  \begin{block}{Conditionals of a multivariate Gaussian}
    Let $\bx=(\bx_A,\bx_B)^T\sim\mathcal{N}(\Vmu,\VSigma)$ with
    $\Vmu=(\Vmu_A,\Vmu_B)^T$ and
    $\VSigma=\bigl(\begin{smallmatrix}\VSigma_A & \VSigma_{AB}\\\VSigma_{AB}^T & \VSigma_B\end{smallmatrix}\bigr)$.
    The conditional $p(\bx_A|\bx_B)$ is Gaussian with:
    \begin{align*}
      \Vmu_{A|B} &= \Vmu_A + \VSigma_{AB}\VSigma_B^{-1}(\bx_B-\Vmu_B)\\
      \VSigma_{A|B} &= \VSigma_A - \VSigma_{AB}\VSigma_B^{-1}\VSigma_{AB}^T
    \end{align*}
    (symmetrically for $p(\bx_B|\bx_A)$).
  \end{block}

  \medskip
  \begin{itemize}
    \item Both marginals and conditionals of a multivariate Gaussian are Gaussian.
    \item Setting $\bx_A=f(\bZ)$, $\bx_B=f(\bX)$ gives the GP predictive distribution.
  \end{itemize}
\end{frame}

%-------------------------------------------------------
\begin{frame}{Posterior Predictive Distribution -- No Noise}{}
  \begin{block}{Predictive distribution (no noise)}
    $p(f(\bZ)\mid\bZ,\bX,f(\bX))=\mathcal{N}(\Vmu_{pr},\,\VSigma_{pr})$ with
    \begin{align*}
      \Vmu_{pr} &= \Vmu_{\bZ}+G_{\bZ,\bX}\,G_{\bX}^{-1}(\btt-\Vmu_{\bX})\\
      \VSigma_{pr} &= G_{\bZ}-G_{\bZ,\bX}\,G_{\bX}^{-1}\,G_{\bZ,\bX}^T
    \end{align*}
  \end{block}

  \medskip
  \begin{itemize}
    \item Under zero-mean prior ($\Vmu_{\bX}=\Vmu_{\bZ}=\bZero$):
      $\Vmu_{pr}=G_{\bZ,\bX}\,G_{\bX}^{-1}\btt$.
    \item Non-zero means are incorporated as a final post-correction step.
    \item The posterior is a \halert{Gaussian}: uncertainty decreases near observed points.
  \end{itemize}
\end{frame}

%-------------------------------------------------------
\begin{frame}{Posterior Samples -- No Noise}{}
  \begin{center}
    \includegraphics[width=0.72\linewidth]{figures/gp_post_sample_no_noise.png}
  \end{center}
  \begin{itemize}
    \item \textcolor{red}{Red}: unknown function $f$; values at 5 training points observed.
    \item \textcolor{blue}{Blue}: samples from the posterior $\mathcal{N}(\Vmu_{pr},\VSigma_{pr})$.
    \item All sampled functions pass \halert{exactly} through the 5 observed points.
  \end{itemize}
\end{frame}

%-------------------------------------------------------
\begin{frame}{Single-Point Prediction -- No Noise}{}
  For a single test point $\bz$:

  \begin{block}{Single-point predictive distribution (no noise)}
    \begin{align*}
      \mu_{pr} &= G_{\bz,\bX}\,G_{\bX}^{-1}(\btt-\Vmu_{\bX})\\
      \sigma^2_{pr} &= \kappa(\bz,\bz)-G_{\bz,\bX}\,G_{\bX}^{-1}\,G_{\bz,\bX}^T
    \end{align*}
  \end{block}

  \medskip
  \begin{itemize}
    \item At each training point $\bx_i$:
          $\mu(f(\bx_i)\mid\bX,\btt)=t_i$ and $\sigma^2=0$.
    \item This is exact \halert{interpolation}: all posterior functions agree at $\bx_i$.
    \item Predictive uncertainty grows away from training points.
  \end{itemize}
\end{frame}

%-------------------------------------------------------
\begin{frame}{Predictive Mean and Variance -- No Noise}{}
  \begin{center}
    \includegraphics[width=0.78\linewidth]{figures/gp_predict_no_noise.png}
  \end{center}
  \begin{itemize}
    \item \textcolor{blue}{Blue}: predictive mean $\mu_{pr}(x)$ over the domain.
    \item Yellow band: $\pm 1$ standard deviation $\sigma_{pr}(x)$.
    \item Variance is \halert{zero} at the 5 observed points (exact interpolation).
  \end{itemize}
\end{frame}

% ============================================================
\section{GP Regression with Gaussian Noise}
% ============================================================

%-------------------------------------------------------
\begin{frame}{Gaussian Noise Model}{Motivation}
  \begin{block}{More realistic assumption}
    Each $t_i$ is a \halert{noisy} observation of $f(\bx_i)$:
    \[
      t_i = f(\bx_i)+\varepsilon,\qquad p(\varepsilon)=\mathcal{N}(0,\sigma_f^2)
    \]
  \end{block}

  \medskip
  \begin{itemize}
    \item Noise adds to the diagonal of the covariance matrix:
      \[
        \VSigma_{\bX}[i,j]=\kappa(\bx_i,\bx_j)\;(i\neq j),
        \qquad
        \VSigma_{\bX}[i,i]=\kappa(\bx_i,\bx_i)+\sigma_f^2
      \]
    \item In matrix form: $\hat{\VSigma}_{\bX}=G_{\bX}+\sigma_f^2\Ii$.
    \item Observations no longer coincide with exact function values.
  \end{itemize}
\end{frame}

%-------------------------------------------------------
\begin{frame}{Noisy GP -- Modified Covariance Matrix}{}
  \begin{block}{Covariance matrix with Gaussian noise}
    \[
      \hat{\VSigma}_{\bX}=G_{\bX}+\sigma_f^2\Ii=
      \begin{pmatrix}
        \kappa(\bx_1,\bx_1)+\sigma_f^2 & \kappa(\bx_1,\bx_2) & \cdots & \kappa(\bx_1,\bx_m)\\
        \kappa(\bx_2,\bx_1) & \kappa(\bx_2,\bx_2)+\sigma_f^2 & \cdots & \kappa(\bx_2,\bx_m)\\
        \vdots & \vdots & \ddots & \vdots\\
        \kappa(\bx_m,\bx_1) & \cdots & \cdots & \kappa(\bx_m,\bx_m)+\sigma_f^2
      \end{pmatrix}
    \]
  \end{block}

  \medskip
  \begin{itemize}
    \item Only \halert{diagonal} entries are modified (noise is i.i.d.\ and independent).
    \item Off-diagonal entries remain determined solely by $\kappa$.
  \end{itemize}
\end{frame}

%-------------------------------------------------------
\begin{frame}{Noisy GP -- Prior Visualisation}{}
  \begin{center}
    \includegraphics[width=0.72\linewidth]{figures/gp_noise.png}
  \end{center}
  \begin{itemize}
    \item Illustration of the effect of Gaussian noise on the GP prior.
    \item Observations are scattered around the true function value.
  \end{itemize}
\end{frame}

%-------------------------------------------------------
\begin{frame}{Joint Distribution -- Noisy Case}{}
  \begin{itemize}
    \item Training: $\bX,\btt$ with $t_i=f(\bx_i)+\varepsilon$.
    \item Test points: $\bZ=(\bz_1,\ldots,\bz_r)$.
  \end{itemize}

  \medskip
  The joint distribution of $(f(\bX),f(\bZ))$ uses the modified covariance:
  \[
    \hat{\VSigma}_{(\bX,\bZ)}=
    \begin{pmatrix}\hat{\VSigma}_{\bX} & G_{\bZ,\bX}^T \\ G_{\bZ,\bX} & G_{\bZ}\end{pmatrix},
    \qquad
    \hat{\VSigma}_{\bX}=G_{\bX}+\sigma_f^2\Ii
  \]

  \medskip
  \begin{itemize}
    \item Applying Gaussian conditionals with $G_{\bX}$ replaced by $\hat{\VSigma}_{\bX}$ yields the posterior.
  \end{itemize}
\end{frame}

%-------------------------------------------------------
\begin{frame}{Posterior Predictive Distribution -- Noisy Case}{}
  \begin{block}{Noisy predictive distribution}
    $p(f(\bZ)\mid\bZ,\bX,\btt)=\mathcal{N}(\hat{\Vmu}_{pr},\,\hat{\VSigma}_{pr})$ with
    \begin{align*}
      \hat{\Vmu}_{pr} &= \Vmu_{\bZ}+G_{\bZ,\bX}\,\hat{\VSigma}_{\bX}^{-1}(\btt-\Vmu_{\bX})\\
      \hat{\VSigma}_{pr} &= G_{\bZ}-G_{\bZ,\bX}\,\hat{\VSigma}_{\bX}^{-1}\,G_{\bZ,\bX}^T
    \end{align*}
    Under zero-mean prior: $\hat{\Vmu}_{pr}=G_{\bZ,\bX}\,\hat{\VSigma}_{\bX}^{-1}\btt$.
  \end{block}

  \medskip
  \begin{itemize}
    \item Same form as noise-free case with $G_{\bX}\to\hat{\VSigma}_{\bX}$.
    \item Posterior variance is now \halert{positive} even at training points.
    \item Noise introduces \halert{smoothing}: no exact interpolation.
  \end{itemize}
\end{frame}

%-------------------------------------------------------
\begin{frame}{No Noise vs.\ Noisy Predictive -- Comparison}{}
  \begin{columns}[t]
    \column{0.48\textwidth}
      \begin{block}{No noise}
        \begin{itemize}
          \item Inverts $G_{\bX}$.
          \item Posterior variance is \emph{zero} at training points.
          \item Exact \halert{interpolation}.
        \end{itemize}
      \end{block}
    \column{0.48\textwidth}
      \begin{block}{With Gaussian noise}
        \begin{itemize}
          \item Inverts $\hat{\VSigma}_{\bX}=G_{\bX}+\sigma_f^2\Ii$.
          \item Posterior variance is \emph{positive} everywhere.
          \item \halert{Smoothing} regression.
        \end{itemize}
      \end{block}
  \end{columns}
\end{frame}

%-------------------------------------------------------
\begin{frame}{Posterior Samples -- Noisy Observations}{}
  \begin{center}
    \includegraphics[width=0.72\linewidth]{figures/gp_post_sample_noise.png}
  \end{center}
  \begin{itemize}
    \item Sampled functions \emph{no longer} pass exactly through training points.
    \item The posterior \halert{smooths} through the observations.
  \end{itemize}
\end{frame}

%-------------------------------------------------------
\begin{frame}{Single-Point Prediction -- Noisy Case}{}
  For a single test point $\bz$:

  \begin{block}{Noisy single-point prediction}
    \begin{align*}
      \mu_{pr} &= m(\bx)+G_{\bz,\bX}\,\hat{\VSigma}_{\bX}^{-1}(\btt-\Vmu_{\bX})\\
      \sigma^2_{pr} &= \kappa(\bz,\bz)-G_{\bz,\bX}\,\hat{\VSigma}_{\bX}^{-1}\,G_{\bz,\bX}^T
    \end{align*}
  \end{block}

  \medskip
  \begin{itemize}
    \item Predictive mean no longer coincides with $t_i$ at training points.
    \item Predictive variance is positive everywhere, including at $\bx_i$.
    \item Uncertainty is \emph{not} fully resolved at observed locations.
  \end{itemize}
\end{frame}

%-------------------------------------------------------
\begin{frame}{Predictive Mean and Variance -- Noisy Case}{}
  \begin{center}
    \includegraphics[width=0.72\linewidth]{figures/gp_predict_noise.png}
  \end{center}
  \begin{itemize}
    \item \textcolor{blue}{Blue}: predictive mean $\mu_{pr}(x)$.
    \item Band: predictive uncertainty $\pm\sigma_{pr}(x)$.
    \item At training points: residual variance $>0$ -- \halert{no exact interpolation}.
  \end{itemize}
\end{frame}

% ============================================================
\section{Kernel Hyperparameter Estimation}
% ============================================================

%-------------------------------------------------------
\begin{frame}{Estimating Kernel Hyperparameters}{Why it matters}
  \begin{alertblock}{Key observation}
    GP predictive performance depends \emph{entirely} on the suitability of the chosen kernel.
  \end{alertblock}

  \medskip
  \begin{itemize}
    \item For the RBF kernel with noise:
      \[
        \kappa(\bx_i,\bx_j)=\sigma_f^2\,
        e^{-\frac{1}{2}(\bx_i-\bx_j)^T\mathbf{M}(\bx_i-\bx_j)}
        +\sigma_y^2\,\delta_{ij}
      \]
    \item Simplest choice: $\mathbf{M}=l^{-2}\Ii$ (isotropic length scale $l$).
    \item Three hyperparameters: $\sigma_f$ (signal variance), $\sigma_y$ (noise variance), $l$ (length scale).
    \item Different values yield \halert{very different} predictive behaviours.
  \end{itemize}
\end{frame}

%-------------------------------------------------------
\begin{frame}{Effect of Kernel Hyperparameters}{}
  \begin{center}
    \includegraphics[width=0.85\linewidth]{figures/gp_params.jpg}
  \end{center}
  \begin{center}
    \small
    (Figure from K.\ Murphy, \emph{Machine Learning: A Probabilistic Perspective}, p.\ 519,
    with $(l,\sigma_f,\sigma_y)$ equal to $(1,1,0.1)$, $(0.3,1.08,0.00005)$, $(3.0,1.16,0.89)$.)
  \end{center}
\end{frame}

%-------------------------------------------------------
\begin{frame}{Hyperparameter Estimation -- Type-II MLE}{}
  \begin{block}{Marginal likelihood (empirical Bayes)}
    Optimise $\log p(\btt\mid\bX,\boldsymbol{\theta})$
    with respect to $\boldsymbol{\theta}=(l,\sigma_f,\sigma_y)$:
    \[
      \log p(\btt\mid\bX,\boldsymbol{\theta})
      =-\tfrac{1}{2}\btt^T\hat{\VSigma}_{\bX}^{-1}\btt
       -\tfrac{1}{2}\log|\hat{\VSigma}_{\bX}|
       -\tfrac{n}{2}\log 2\pi
    \]
  \end{block}

  \medskip
  \begin{itemize}
    \item The marginal likelihood automatically balances data-fit and model complexity.
    \item Gradient-based optimisation (e.g.\ L-BFGS) is typically used.
    \item Multiple random restarts recommended to avoid local optima.
    \item After optimisation, $\boldsymbol{\theta}^*$ is used for prediction.
  \end{itemize}
\end{frame}

% ============================================================
\section{Summary}
% ============================================================

%-------------------------------------------------------
\begin{frame}{Summary}{Nonparametric regression at a glance}
  \begin{columns}[t]
    \column{0.48\textwidth}
      \begin{block}{Kernel / local methods}
        \begin{itemize}
          \item \textbf{Equiv.\ kernel}: prediction $=$ weighted sum of $t_i$.
          \item \textbf{Nadaraya-Watson}: normalised kernel weights.
          \item \textbf{LOESS}: local weighted least squares.
          \item \textbf{Local logistic}: weighted cross-entropy for classification.
        \end{itemize}
      \end{block}
    \column{0.48\textwidth}
      \begin{block}{Gaussian processes}
        \begin{itemize}
          \item Prior $\mathcal{GP}(m,\kappa)$ over functions.
          \item Posterior $p(f\mid\bX,\btt)$ is Gaussian.
          \item \textbf{No noise}: exact interpolation.
          \item \textbf{With noise}: smoothing regression.
          \item Hyperparameters: marginal likelihood.
        \end{itemize}
      \end{block}
  \end{columns}
\end{frame}

%-------------------------------------------------------
\begin{frame}{Key Formulas at a Glance}{}
  \begin{center}
    \small
    \renewcommand{\arraystretch}{1.45}
    \begin{tabular}{ll}
      \hline\hline
      \textbf{Result} & \textbf{Formula} \\
      \hline
      Equivalent kernel & $\kappa(\bx_1,\bx_2)=\beta\,\ox_1^T\bS\ox_2$ \\
      Prediction (equiv.\ kernel) & $h(\bx)=\sum_i\kappa(\bx,\bx_i)\,t_i$ \\
      Nadaraya-Watson weights & $w_i(\bx)=\kappa_h(\bx-\bx_i)/\sum_j\kappa_h(\bx-\bx_j)$ \\
      LOESS solution & $\ow^*(\bx)=(\oX^T\bPsi\oX)^{-1}\oX^T\bPsi\,\btt$ \\
      GP prior & $p(f(\bX))=\mathcal{N}(\Vmu_{\bX},\,G_{\bX})$ \\
      GP post.\ mean (no noise) & $\Vmu_{pr}=G_{\bZ,\bX}\,G_{\bX}^{-1}\btt$ \\
      GP post.\ mean (noise) & $\hat{\Vmu}_{pr}=G_{\bZ,\bX}\,\hat{\VSigma}_{\bX}^{-1}\btt$ \\
      GP post.\ cov.\ (no noise) & $\VSigma_{pr}=G_{\bZ}-G_{\bZ,\bX}G_{\bX}^{-1}G_{\bZ,\bX}^T$ \\
      GP post.\ cov.\ (noise) & $\hat{\VSigma}_{pr}=G_{\bZ}-G_{\bZ,\bX}\hat{\VSigma}_{\bX}^{-1}G_{\bZ,\bX}^T$ \\
      \hline\hline
    \end{tabular}
  \end{center}
\end{frame}

%\end{document}
